\scanpage{14}
\[
(18)\qquad\widehat f:\overset{\deleted{\wedge}}{X}\mathbin{-}Y
\]
factors as
\diagram{\begin{tikzcd}[ampersand replacement=\&]
(19)\quad X\arrow[r,"\widehat f_0","\sim"']\&Y'\arrow[r,"u"]\&Y
\end{tikzcd}}
where $u$ is a $C_0$-morphism, and $\widehat f_0$ a $\widehat C_0$-\uline{isomorphism}
(\uline{NB} $Y'$ is an object $C_0$-isomorphic to $X$), and where $u$ is
a morphism of $C_0$. In the case where one
assumes that for every arrow $u:Y'\mathbin{-}Y$ in $C$,
the composite homomorphism
\diagram{\begin{tikzcd}[ampersand replacement=\&]
\operatorname{Aut}_Y(Y')\arrow[r]\&\operatorname{Aut}(Y')\arrow[r,no head]\&\operatorname{Aut}(Y')^{\wedge}
\end{tikzcd}}
\uline{injective}, one sees that $Y'$ is determined
in terms of $\widehat f$ up to unique $Y$-isomorphism. \deleted{In} The
\inserted{recipe} \deleted{addition, for} for composing two morphisms defined
thus, $\widehat f=u\widehat f_0$ and $\widehat g=v\widehat g_0$
\margin{\uline{NB} The canonical functor $C_0\longrightarrow\widehat C_0$ is conservative (i.e. $\widehat f_0$ is a $C_0$-isomorphism $\Longleftrightarrow$ it is a $\widehat C_0$-isomorphism) \deleted{In \unsure{addition}} \inserted{moreover,} Every $\widehat C_0$-isomorphism is a limit of $C_0$-isomorphisms\ldots\ in a precise way,
$\operatorname{Isom}_{\widehat C_0}(X,Y)\simeq\operatorname{Isom}_{C_0}(X,Y)\mathbin{\wedge}_{\operatorname{Aut}_{C_0}(X)}\widehat{\operatorname{Aut}}_{C_0}(X)$
and $\operatorname{Aut}_{\widehat C_0}(X)\simeq\widehat{\operatorname{Aut}}_{C_0}(X)$.}
\diagram{\begin{tikzcd}[ampersand replacement=\&]
(20)\quad Y\arrow[r,no head,"\deleted{\unsure{g}}\widehat g_0","\sim"']\&Z'\arrow[r,no head,"v"]\&Z,
\end{tikzcd}}
and putting them into the same standard form,
is the following: one approximates $\widehat g_0:Y\overset{\sim}{\longrightarrow}Z'$
(which is a $\widehat C_0$-isomorphism) by
a \deleted{genuine} \inserted{$C_0$-isomorphism} morphism \deleted{$\unsure{u}$} $w$
in such a way that one has
\[
\widehat g_0=w\widehat\varphi\qquad\text{with}\quad\widehat\varphi\in\widehat H_u\subset\widehat{\operatorname{Aut}}_{C_0}(Y)
\]
\diagram{\begin{tikzcd}[ampersand replacement=\&,column sep=small]
\&\&Y'\arrow[r,"u"]\&Y\arrow[dr,"w"]\&\&\\
(21)\quad X\arrow[r,no head,"\widehat f_0","\sim"']\&Y'\arrow[r,"u"]\arrow[ur,"\widehat\psi?" ]\&Y\arrow[rr,"\widehat g_0","\sim"']\arrow[ur,"\widehat\varphi?"]\&\&Z'\arrow[r,"v"]\&Z
\end{tikzcd}}
\margin{\uline{NB} One must know how to interpret in terms of $C_0$ already \figure{\begin{tikzpicture}[baseline=0pt]\draw (0,0)--(.35,0);\draw (0,.08)--(.35,.08);\end{tikzpicture}} the relation $\widehat\varphi u=u\widehat\psi$ as meaning $(\widehat\psi,\widehat\varphi)\in\widehat G_u$, where $G_u\subset\operatorname{Aut}(X)\times\operatorname{Aut}(Y)$ is the stabilizer of $u$\ldots}
(which is possible, since $\widehat H_u$ is open in $\widehat{\operatorname{Aut}}_{C_0}(Y)$),
and one chooses $\widehat\psi\in\widehat{\operatorname{Aut}}(Y')$ such that one has $\widehat\varphi u=u\widehat\psi$
so one will have
\[
(22)\qquad\widehat g\widehat f=(v\widehat g_0)(u\widehat f_0)
=(v\deleted{\unsure{g}}w\widehat\varphi)(u\widehat f_0)
=\underbrace{(vwu)}_{\in\mathrm{Fl}\,C_0}\underset{\in\mathrm{Flisom.}(\widehat C_0)}{\underline{(\widehat\psi\widehat f_0)}}
\]
\margin{13}
